% Werke X/1 (1917), native220–223, printed213–216.
\begingroup
\makeatletter\tagsleft@true\let\veqno\@@leqno\makeatother
\clearpage\SourceLeaf{220}\NoSourceHead
\SourceGroup{VIII}
\HistoricTitle{[DER BILINEARE ALGORITHMUS.]}
\HistoricCitation{[Aus Scheda An, Cereri, Palladi, Iunoni sacrum, Febr.\ 1805.]}
\NumberedItem{1}
\noindent [S.\ 37]
\par\indent Es sind gegeben vier Grössen \(a,b,A,B\); man bildet daraus nach folgender Ordnung
{\small\setlength{\arraycolsep}{3pt}
\[\begin{array}{lllll}
2a'=a+b,&b'b'=ab,&cc=aa-bb,&c=\frac12({}'a-{}'b),&cc=4a'c',\\[.5em]
2a''=a'+b',&b''b''=a'b',&c'^{2}=a'^{2}-b'^{2},&c'=\frac12(a-b),&c'^{2}=4a''c'',\\[.5em]
2a'''=a''+b'',&b'''^{2}=a''b'',&c''^{2}=a''^{2}-b''^{2},&c''=\frac12(a'-b'),&c''^{2}=4a'''c'''
\end{array}\]
\begin{center}etc.\end{center}}
\[\text{Grenze}=h.\]
\par\indent Man mache
\begin{equation*}\tag*{[1]}
\left\{\begin{aligned}
x^2&=\frac{c}{4a'^{\frac12}a''^{\frac14}a'''^{\frac18}\ldots}
=\frac{ca'^{\frac12}a''^{\frac14}a'''^{\frac18}\ldots}{4a'a''^{\frac12}a'''^{\frac14}\ldots}\\[.8em]
&=\frac{c}{4a'}\left(\frac{a'}{a''}\right)^{\frac12}\left(\frac{a''}{a'''}\right)^{\frac14}\left(\frac{a'''}{a^{\mathrm{IV}}}\right)^{\frac18}\ldots\\[.8em]
&=\frac{cb'^{\frac12}b''^{\frac14}\ldots}{4b''b'''^{\frac12}\ldots}
\end{aligned}\right.
\end{equation*}
\[x^4=\frac{a-b}{8a''}\left(\frac{a''}{a'''}\right)^{\frac12}\left(\frac{a'''}{a^{\mathrm{IV}}}\right)^{\frac14}\text{etc.,}\]
\noindent so ist
\begin{equation*}\tag*{[2]}
\left\{\begin{aligned}
a&=h(1+2x^4+2x^{16}+2x^{36}+\ldots)^2\\
b&=h(1-2x^4+2x^{16}-2x^{36}+\ldots)^2\\
c&=h(2x+2x^9+2x^{25}+\ldots)^2.
\end{aligned}\right.
\end{equation*}
\clearpage\SourceLeaf{221}\ArithmeticHead{214}{L}{ANALYSIS.\quad NACHLASS.}
\noindent Ferner
\begin{equation*}\tag*{[3]}
\left\{\begin{gathered}
\begin{array}{ll}
A'=\frac12(A+B),& B'=\dfrac{2ABa'}{b'(A+B)}\\[.7em]
A''=\frac12(A'+B'),& B''=\dfrac{2A'B'a''}{b''(A'+B')}\\[.7em]
A'''=\frac12(A''+B''),& B'''=\dfrac{2A''B''a'''}{b'''(A''+B'')}
\end{array}\\
\text{etc.}
\end{gathered}\right.
\end{equation*}
\noindent Man bestimme \(H\) durch die Gleichung
\begin{equation*}\tag*{[4]}
H=\frac{hA'^{\frac12}A''^{\frac14}A'''^{\frac18}\ldots}{a'^{\frac12}a''^{\frac14}a'''^{\frac18}\ldots}
=\frac{chA'^{\frac12}A''^{\frac14}A'''^{\frac18}\ldots}{4x}
\end{equation*}
\noindent und \(y\) durch [die weiter unten folgende Gleichung [7]], so ist
\begin{equation*}\tag*{[5]}
\left\{\begin{aligned}
A&=H\left(1+\left(y+\frac1y\right)x^4+\left(yy+\frac1{yy}\right)x^{16}+\text{etc.}\right)^2\\[.5em]
B&=H\left(1-\left(y+\frac1y\right)x^4+\left(yy+\frac1{yy}\right)x^{16}-\text{etc.}\right)^2
\end{aligned}\right.
\end{equation*}
\NumberedItem{2}
\noindent [S.\ 38]
\[\begin{gathered}
\frac ba=\cos2M,\quad\frac{a'}a=\cos M^2\\[.4em]
\frac{b'}{a'}=\cos2M'\\
\text{etc.}
\end{gathered}\]
\[\operatorname{tg}M^2=\sin2M'\]
\[a^\infty=a\cos M^2\cos M'^{2}\cos M''^{2}\text{ etc.}\]
\begin{equation*}\tag*{[6]}
\left\{\begin{gathered}
\sin\varphi\;\surd\sin2M=\operatorname{tg}\psi\\
\sin\varphi'\;\surd\sin2M'=\operatorname{tg}\psi'\\
\sin\varphi''\;\surd\sin2M''=\operatorname{tg}\psi''\\
\text{etc.}\\
\sin2\psi'=\operatorname{tg}\psi\;\surd\sin2M
\end{gathered}\right.
\end{equation*}
\[\frac{\sin\varphi'}{\sin\varphi}=\frac{\cos M^2}{\cos\psi'^{2}}\]
\clearpage\SourceLeaf{222}\ArithmeticHead{215}{R}{DER BILINEARE ALGORITHMUS.}
{\small
\[\begin{aligned}
\sin\varphi^\infty&=\sin\varphi\frac{\cos M^2.\cos M'^{2}.\cos M''^{2}\ldots}{\cos\psi'^{2}.\cos\psi''^{2}.\cos\psi'''^{2}\ldots}\\
&=\frac{a^\infty\sin\varphi}{a\cos\psi'^{2}.\cos\psi''^{2}.\cos\psi'''^{2}\ldots}
\end{aligned}\]}
\noindent [Es folgt ein Beispiel \(2M=75^0\)]
\ClosingRule
\par\medskip\noindent [S.\ 39]
\[\begin{aligned}
xx&=\frac{\sin2M}{4\cos M^2.\cos M'.\cos M''^{\frac12}.\cos M'''^{\frac14}\ldots}\\[.6em]
&=\frac{\operatorname{tang}M}{2\cos M'.\cos M''^{\frac12}.\cos M'''^{\frac14}\ldots}
\end{aligned}\]
\[\frac{2x+2x^9+\ldots}{1+2x^4+\mathord{.}\;\mathord{.}}=\surd\sin2M\]
\[\frac{1-2x^4+2x^{16}+\ldots}{1+2x^4+2x^{16}+\ldots}=\surd\cos2M\]
\[\frac{\left(y+\frac1y\right)x+\left(y^3+\frac1{y^3}\right)x^9+\ldots}{1+\left(yy+\frac1{yy}\right)x^4+\left(y^4+\frac1{y^4}\right)x^{16}+\ldots}=\operatorname{tg}\psi\]
\begin{equation*}\tag*{[7]}y+\frac1y=2\sin\varphi^{(\infty)}\end{equation*}
\noindent [Es folgen Zahlenrechnungen].
\NumberedItem{3}
\noindent [S.\ 40]
\par\indent Eine andere Manier:
\[\begin{aligned}
\operatorname{tg}\omega&=\operatorname{tg}\varphi\;\surd(1-\sin2M^2.\sin\varphi^2)\\[.6em]
\operatorname{tg}\omega'&=\operatorname{tg}2\omega.\surd\cos2M=\operatorname{tg}2\omega.\surd\frac ba\\[.6em]
\operatorname{tg}\omega''&=\operatorname{tg}2\omega'.\surd\cos2M'=\operatorname{tg}2\omega'.\surd\frac{b'}{a'}\\[.6em]
\operatorname{tg}\omega'''&=\operatorname{tg}2\omega''.\surd\cos2M''=\operatorname{tg}2\omega''.\surd\frac{b''}{a''}
\end{aligned}\]
\begin{center}etc.\end{center}
\noindent Dann ist \(\varphi^\infty\) das letzte Glied der Reihe \(\omega,\;\frac12\omega',\;\frac14\omega'',\;\frac18\omega'''\) etc.
\clearpage\SourceLeaf{223}\ArithmeticHead{216}{L}{ANALYSIS.\quad NACHLASS.}
{\small
\begin{equation*}\tag*{[8]}
\sin\varphi=\frac{1+2x^4+\ldots}{2x+2x^9+\ldots}.
\frac{\left(y+\frac1y\right)x+\left(y^3+\frac1{y^3}\right)x^9+\ldots}{1+\left(yy+\frac1{yy}\right)x^4+\ldots}
\end{equation*}
\[\begin{gathered}
\surd(1-\sin2M^2.\sin\varphi^2)=\\[-.5em]
\frac{\surd\,1-2xx+2x^8+\ldots)\;\surd\left(1-xx\left(yy+\frac1{yy}\right)+x^8\left(y^4+\frac1{y^4}\right)+\ldots\right)}{1+\left(yy+\frac1{yy}\right)x^4+\ldots}
\end{gathered}\]}
\noindent [Es folgen Umformungen der hier auftretenden Reihen in Produkte, die in der Abhandlung Werke III, S.\ 446 ff.\ viel ausführlicher und übersichtlicher dargestellt sind, und die darum nicht abgedruckt werden; auf derselben Seite steht noch unten die Formel:]
{\small
\begin{equation*}\tag*{[9]}
\begin{aligned}
\operatorname{tang}\varphi={}&\left[\frac1i.\right]
\frac{\left(y+\frac1y\right)x+\left(y^3+\frac1{y^3}\right)x^9+\ldots}{\left(y-\frac1y\right)x-\left(y^3-\frac1{y^3}\right)x^9+\ldots}\\[.6em]
&.\frac{1+2x^4+2x^{16}+\ldots}{1-2x^4+2x^{16}-\ldots}\;\text{[*)]}
\end{aligned}
\end{equation*}}
\par\medskip\noindent\rule{.25\linewidth}{.4pt}
\par\smallskip\noindent [*) In der Handschrift lautet der zweite Faktor
\[\frac{1-2x^4+2x^{16}\ldots}{1+2x^4-2x^{16}\ldots}\mathpunct{\raisebox{.5ex}{.}}\Bigr]\]
\ClosingRule
\clearpage\endgroup
